\documentclass[lualatex,ja=standard]{bxjsarticle}
\usepackage{amsmath,amssymb}
\usepackage{enumitem}
\setlist[enumerate]{itemsep=0.6em, topsep=0.5em}
\pagestyle{empty}

\begin{document}

\noindent\textbf{数IA 共通テスト演習（図形・確率・データ分析）}
\hfill\textbf{制限時間：25分}
\hrule
\vspace{0.8em}

\begin{enumerate}
\item 三角形 $ABC$ において $AB=6,\ AC=8,\ \angle A=120^\circ$ とする。
\begin{enumerate}[label=(\arabic*)]
\item $BC$ の長さを求めよ。
\item 三角形 $ABC$ の面積を求めよ。
\end{enumerate}

\item 三角形 $ABC$ で $AB=5,\ BC=7,\ CA=8$ とする。
$\cos B$ の値を求め、$\angle B$ が鋭角か鈍角か判定せよ。

\item 1から9までの整数が1つずつ書かれたカード9枚から、同時に2枚を取り出す。
取り出した2枚の数の和が6の倍数である確率を求めよ。

\item 赤玉4個、白玉3個、青玉2個が入った袋から、同時に3個取り出す。
「3色のうちちょうど2色が出る」確率を求めよ。

\item 2個のさいころを同時に投げる。
和が9以上であることがわかっているとき、積が偶数である条件付き確率を求めよ。

\item ある高校のA組40人の数学テストについて、平均点は64点、中央値は66点であった。
次の記述のうち、必ず正しいものをすべて選べ。
\begin{itemize}
\item[(ア)] 66点以上の生徒は20人以上いる。
\item[(イ)] 64点以上の生徒は20人以上いる。
\item[(ウ)] 最高点は64点以上である。
\end{itemize}

\item 次のデータ
\[
4,\ 5,\ 6,\ 7,\ 7,\ 8,\ 9,\ 10,\ 12,\ 12
\]
について、四分位数 $Q_1,\ Q_2,\ Q_3$ と四分位範囲を求めよ。

\item ある模試の数学の得点が平均58点、標準偏差12点のとき、得点74点の受験者の偏差値を求めよ。
また、偏差値65に相当する得点を求めよ。
\end{enumerate}

\end{document}
